\section*{Erd\H{o}s Problem \#364: three consecutive powerful integers}
\subsection*{1. Statement and scope}
A positive integer is powerful, or $2$-full, when every prime exponent
in it is either zero or at least two. Are there any three consecutive
positive powerful integers? We prove structural restrictions,
reconstruct the infinite family of consecutive pairs, and give the
precise conditional consequence of $abc$. An exact finite search is
included as a bounded check, not as an answer to the infinite question,
which remains unresolved.

\subsection*{2. Unique parametrization and an exhaustive search}
Every powerful integer has a unique expression
\[
 n=a^2b^3,\qquad a,b\ge1,\quad b\text{ squarefree}.
\]
For each odd positive prime exponent $e$, place that prime once in
$b$ and $(e-3)/2$ times in $a$; for an even exponent use $e/2$ copies
in $a$. This construction is forced by parity and proves uniqueness.
There is no requirement that $a$ and $b$ be coprime.
In particular the number of powerful integers at most $X$ is at most
\[
 \sum_{b\le X^{1/3}}\left\lfloor\sqrt{X/b^3}\right\rfloor
 \le \sqrt X\sum_{b\ge1}b^{-3/2}.
\]
Enumerating precisely the squarefree $b\le X^{1/3}$ and
$1\le a\le\lfloor\sqrt{X/b^3}\rfloor$ therefore gives an exhaustive
finite algorithm, with no factorization heuristic.
For $X=10^{10}$ the exact enumeration contains $214122$ integers.
Checking membership of $n,n+1,n+2$ finds no triple all of whose
members lie in $[1,X]$. This computation makes no claim about larger
integers and is far smaller than published computational searches.

\subsection*{3. Local obstructions and an infinite family of pairs}
Four consecutive integers cannot all be powerful, because one has
residue $2$ modulo $4$. If $m-1,m,m+1$ are powerful, this also shows
that the center satisfies $m\equiv0\pmod4$.
Residues $3,6$ modulo $9$ are impossible for a powerful integer,
so examining the three consecutive residues gives
\[
 m\equiv0,1,8\pmod9.
\]
The two odd neighbors cannot both be squares: two positive squares
differing by $2$ would give $(v-u)(v+u)=2$, which is impossible for
positive integers $u<v$. Thus at least one neighbor has a nontrivial
squarefree cube factor in the parametrization above.

These restrictions do not prevent consecutive \emph{pairs}.
Starting at $(x_1,y_1)=(3,1)$, define
\[
 x_{j+1}=3x_j+8y_j,\qquad y_{j+1}=x_j+3y_j.
\]
Direct expansion preserves $x_j^2-8y_j^2=1$, and both coordinates
increase. Hence $8y_j^2$ and $8y_j^2+1=x_j^2$ give infinitely many
distinct powerful pairs. Every prime exponent in the first number
is even except possibly that of $2$, which is $3+2v_2(y_j)$ and
therefore at least three.

\subsection*{4. What the $abc$ conjecture would actually imply}
Assume the usual $abc$ conjecture: for every $\varepsilon>0$ there is
$C_\varepsilon$ such that coprime positive $a+b=c$ satisfy
$c\le C_\varepsilon\operatorname{rad}(abc)^{1+\varepsilon}$.
If $m-1,m,m+1$ are powerful, each obeys
$\operatorname{rad}(u)\le\sqrt u$. Apply the conjecture to
\[
 (m^2-1)+1=m^2.
\]
The summands are coprime, and
\[
 \operatorname{rad}\big((m^2-1)m^2\big)
 =\operatorname{rad}\big((m-1)m(m+1)\big)
 \le\sqrt{(m-1)m(m+1)}<m^{3/2}.
\]
It follows that $m^2\le C_\varepsilon m^{3(1+\varepsilon)/2}$.
Taking $\varepsilon<1/3$ bounds $m$.
Thus $abc$ implies only finitely many such triples, not that there
are none. No use of this conjecture is made in the unconditional
results or the finite search.

\subsection*{5. Assessment and references}
The remaining obligation is an unconditional existence proof or an
obstruction valid for every candidate center. The congruences,
parametrization, and conditional finiteness do not provide it.
Source: \texttt{https://www.erdosproblems.com/364}, checked 2026-09-25,
which records Mahler's Pell observation and the conditional $abc$
consequence. The prior GPT Pro 5.2 argument was checked; its
parametrization and elementary obstruction are reconstructed here.
No novelty or formal Lean verification is claimed.
\subsection*{COMPLETION ESTIMATE}
\textbf{COMPLETION: 35\%}
